% Sections 18.1-18.8, from lectures/L18/appendix/a_register_bank.md.

\section{Why this layer exists}
\label{c18:sec:why}

\code{can_controller}'s status outputs are made for a caller in the same clock domain that sees every
cycle: \code{tx_done} and \code{rx_valid} are \textbf{one-cycle pulses}, 20~ns at 50~MHz
(\chapref{c11:ch}). A driver polling over SPI samples the system a few \emph{microseconds} at a time
at best \textemdash{} a pulse is gone thousands of cycles before the next poll arrives. The register
map (appendix~\ref{app:protocol}) therefore promises something else: \textbf{sticky, pollable levels}
with explicit clears (\code{RX_ACK}, \code{ERROR_FLAGS}), and write-triggered events
(\code{TX_SEND}).

Converting between the two worlds is this module's entire job. It contains no protocol knowledge in
either direction: nothing about CAN frames (that stays in \code{can_controller}) and nothing about
SPI (that comes later in this chapter, and in \chapref{c19:ch}). It is registers, latches and a pulse
generator \textemdash{} which is exactly why it can be verified completely during a single session,
before any SPI exists at all.

You have seen both of its core tricks before, in \chapref{c3:ch} and \ref{c4:ch}. A sticky status bit
is \chapref{c3:ch}'s flip-flop with the two questions answered explicitly: what sets it, and what
clears it. And turning a held level into a one-cycle pulse is exactly what \code{button_sync}'s edge
detection did in \chapref{c4:ch} \textemdash{} here it runs the same way for \code{TX_SEND} (a write
becomes a pulse) and the opposite way for \code{tx_done} (a pulse becomes a level that stays). This
chapter promotes both from an exercise to a real architectural layer with a documented contract.

% ------------------------------------------------------------------------------------------
\section{The interface}
\label{c18:sec:iface}

Sixteen ports, \textbf{every input first, then every output}, the same convention as every module in
the project \textemdash{} and, as everywhere else in it, the binding is \emph{positional}, so the
order below is the contract \code{register_bank_tb} holds you to:

\begin{textdrawing}
 1 clock        5 reg_wdata    9 rx_data     13 tx_id
 2 reset_s2_n   6 tx_done     10 rx_valid    14 tx_dlc
 3 reg_index    7 rx_id       11 error       15 tx_data
 4 reg_write    8 rx_dlc      12 reg_rdata   16 tx_req
\end{textdrawing}

The types come from \code{can_def} (\code{use work.can_def.all;}), so the ports facing the controller
have exactly the same types as \code{can_controller}'s own. Besides \code{clock} and
\code{reset_s2_n} (ports 1 and 2), the other fourteen fall into two groups.

\paragraph{The side facing the bridge (ports 3--5 and 12)}
A deliberately minimal register access port \textemdash{} what \chapref{c19:ch}'s bridge will drive,
and what the testbench drives directly today:
\begin{itemize}
  \item \code{reg_index : std_logic_vector(3 downto 0)} \textemdash{} which register, the protocol's
    index (offset divided by 4).
  \item \code{reg_write : std_logic} \textemdash{} a \textbf{one-cycle strobe}; at its rising clock
    edge \code{reg_wdata} is committed to the addressed register.
  \item \code{reg_wdata : std_logic_vector(31 downto 0)} \textemdash{} the value to write.
  \item \code{reg_rdata : std_logic_vector(31 downto 0)} \textemdash{} the addressed register's
    current value, \textbf{combinationally}: it follows \code{reg_index} with no clock edge involved.
    The latch-once-on-read rule from the protocol specification is the \emph{bridge's} job
    (\chapref{c19:ch}), not this module's \textemdash{} the bank always tells the current truth, and
    the bridge decides when it takes its single atomic sample.
\end{itemize}

\paragraph{The side facing the controller (ports 6--11 and 13--16)}
Mirrors \code{can_controller}'s register-side ports one for one: \code{tx_id}, \code{tx_dlc},
\code{tx_data} and \code{tx_req} out to the controller, and \code{tx_done}, \code{rx_id},
\code{rx_dlc}, \code{rx_data}, \code{rx_valid} and \code{error} back from it. In \chapref{c19:ch}'s
top level they are wired straight across, port to port, with nothing in between.

\code{reset_s2_n} is the \emph{synchronised} reset, as for every subblock in the project.
\Chapref{c19:ch}'s top level owns a \code{meta_prev} for it, exactly as \code{can_controller} owns its
own internally.

% ------------------------------------------------------------------------------------------
\section{The three STATUS bits: one latch each}
\label{c18:sec:status}

Every \code{STATUS} bit is a flip-flop with a named set condition and a named clear condition, and
nothing more:

\begin{booktable}{@{}p{2.2em}p{5em}L>{\raggedright\arraybackslash}p{9em}@{}}
  \thead{Bit} & \thead{Meaning} & \thead{Set by} & \thead{Cleared by} \\
  \midrule
  0 & TX ready & reset; the controller's \code{tx_done} pulse; \textbf{a rising edge on \code{error}
    while this bit is low}; \textbf{an accepted write to \code{TX_ABORT}} & an accepted write to
    \code{TX_SEND} \\
  1 & RX valid & the controller's \code{rx_valid} pulse & a write of \code{0x1} to
    \code{RX_ACK} \\
  2 & Error & a \textbf{rising edge} on the controller's \code{error} level & a write of \code{0x0}
    to \code{ERROR_FLAGS} \\
\end{booktable}

Three of them deserve a second look:
\begin{itemize}
  \item \textbf{Bit 0 starts at \code{'1'}.} Out of reset the controller is idle, so the bank is ready
    to take a frame \textemdash{} a driver polling before it has ever transmitted has to see
    ``ready'', or \code{send()} would fail forever on a freshly powered system.
  \item \textbf{Bit 2 latches an edge, not a level.} The controller holds \code{error} high until the
    next frame begins (\chapref{c11:ch}), on its own schedule. Latching the rising edge gives the
    \emph{driver} a life cycle of its own for the flag: a clear removes it even while the
    controller's level is still high, and the latch is not set again until a genuinely \emph{new}
    error produces a new rising edge. Latching the level instead would make the clear appear broken
    \textemdash{} cleared and immediately set again \textemdash{} for as long as the controller holds
    the line.
  \item \textbf{Bit 0 has three ways back up, and that is deliberate.} Read the two extra ones as belt
    and braces around the one failure that can take the whole node out.
\end{itemize}

\begin{aside}
A transmission losing arbitration ends without ever reaching EOF. \Chapref{c17:ch} has the controller
pulse \code{tx_done} on that path too, precisely so that this bit comes back \textemdash{} but the
controller is \emph{yours}, and a version that forgets leaves the bit low forever, so every later
\code{TX_SEND} is discarded and the node cannot transmit again until it is reset. The bank therefore
does not rely on the controller getting it right. \textbf{A rising edge on \code{error} while bit 0 is
already low can only mean that a transmission ended badly}, since a node is either transmitting or
receiving, and a reception error cannot occur while a transmission is in progress. So that edge sets
the bit too, and no \code{role} port is needed to know it.

\code{TX_ABORT} is the third path: an explicit escape hatch the driver can pull if it judges that it
has waited long enough. Nothing in the design should ever need it. Build it in anyway \textemdash{}
the cost is a single write-only register, and the alternative is a node a driver has no way of
freeing.
\end{aside}

A single ordering rule ties the latches together: in the clocked process, \textbf{the controller's
events are applied after the bridge's writes}, so that if a \code{TX_SEND} write and a
\code{tx_done} pulse ever landed on the same edge the hardware wins and the bank stays ready. A
conflict that rare should be resolved in the direction that cannot deadlock the system, and that is
the same principle the paragraph above applies three times over.

% ------------------------------------------------------------------------------------------
\section{\code{TX_SEND}: a write becomes a pulse}
\label{c18:sec:txsend}

A \code{TX_SEND} write with \code{reg_wdata(0) = '1'}, arriving while TX ready is set, does exactly
two things at that clock edge: raises \code{tx_req} \textbf{for one cycle}, and clears TX ready. More
precisely: \code{tx_req} is high for the single cycle \emph{after} the write's clock edge, and low
again at the next one \textemdash{} the testbench checks both ends of that claim.

Everything else about this register is what it \emph{ignores}:
\begin{itemize}
  \item A write while TX ready is low (a transmission is in progress) is discarded entirely
    \textemdash{} no pulse, no state change. Overwriting \code{tx_id}, \code{tx_dlc} or
    \code{tx_data} mid-transmission is the corruption the driver's own ready check guards against
    from the software side; this is the same guard on the hardware side, and neither side may assume
    the other is there.
  \item A write with \code{reg_wdata(0) = '0'} does nothing: the register map says ``write \code{0x1}
    to trigger'', and the bank holds it to that.
  \item Reads give zero \textemdash{} \code{TX_SEND} stores nothing, it \emph{does} something.
\end{itemize}

Why the pulse discipline matters shows most clearly in the failure it prevents. \code{can_controller}
acts on \code{tx_req} as soon as it is high and the bus is idle (\chapref{c16:ch}), so a \code{tx_req}
left \emph{held} requests afresh the instant a frame ends, forever, wiping \code{error} at every pass
back through the start state. The request has to be an event, and something has to make it one. At
this layer that is easier \textemdash{} a write strobe is already an event, not a level \textemdash{}
but it still has to be \emph{enforced}: one write, one pulse, whatever the master does next.

% ------------------------------------------------------------------------------------------
\section{\code{TX_ABORT}: the escape hatch}
\label{c18:sec:txabort}

Index 12, write-only, and the mirror image of \code{TX_SEND}: a write of \code{0x1} sets TX ready back
to \code{'1'}. A write of anything else does nothing, and a read gives \code{0x00000000}.

It touches nothing else. It does not reach \code{can_controller} \textemdash{} the controller has no
abort input, its fifteen ports were fixed in \chapref{c11:ch} and \code{can_controller_tb} binds them
positionally \textemdash{} and it needs none. A controller that has aborted a transmission is already
back in \code{STATE_IDLE} and perfectly able to transmit again; the only thing stuck is this bank's
sticky bit. \code{TX_ABORT} frees it, and that is the whole module.

It is worth a moment's pause, since it is the bug's general form. Nothing in the \emph{hardware} had
deadlocked. What had deadlocked was a \textbf{status bit whose only set condition hung on an event
that no longer occurred}. Every sticky flag with one way up and one way down is one missed event away
from the same failure, and the fix is always the same: make sure every exit from the busy state also
raises the flag, and give the caller a way to force it in case you are wrong about ``every''.

% ------------------------------------------------------------------------------------------
\section{The RX side: capturing at the pulse}
\label{c18:sec:rx}

When \code{rx_valid} pulses, the bank does two things at the same edge: sets STATUS bit 1, and
\textbf{captures} \code{rx_id}, \code{rx_dlc} and \code{rx_data} into the four \code{RX_*} registers
\textemdash{} \code{rx_data}'s upper word (bits 63--32, data bytes 0--3, byte 0 in the most
significant position) into \code{RX_DATA_LO}, its lower word into \code{RX_DATA_HI}.

Capturing, instead of wiring the controller's ports straight to the read multiplexer, is what gives
the driver a stable frame to read: the \code{RX_*} registers hold their values while the driver works
its way through four separate SPI read transactions, whatever the controller's ports do meanwhile. A
write to \code{RX_ACK} clears \emph{only the flag} \textemdash{} the data registers keep their
contents until the next capture, which is harmless: the register map's contract is that \code{RX_*} is
meaningful \emph{while RX valid is set}.

What the capture does \textbf{not} give the driver is buffering. A new frame completing before the
previous one has been acknowledged is simply captured over it \textemdash{} one frame is held, the
newest wins. That is not this module being lazy; it is a documented gap inherited from the controller
itself, which holds exactly one received frame and overwrites it at the next (\chapref{c11:ch}). The
bank makes that gap visible at the register level instead of hiding it, and \chapref{c19:ch} returns
to what it costs a driver.

% ------------------------------------------------------------------------------------------
\section{Masking on write}
\label{c18:sec:mask}

Writable registers store only the bits that exist in hardware, and read back as zeros above them:

\begin{booktable}{@{}p{9em}p{9em}L@{}}
  \thead{Register} & \thead{Bits stored} & \thead{A write of \code{0xFFFFFFFF} reads back as} \\
  \midrule
  \code{TX_ID} & 10--0 (\code{id_t}) & \code{0x000007FF} \\
  \code{TX_DLC} & 3--0 (\code{dlc_t}) & \code{0x0000000F} \\
  \code{TX_DATA_LO}/\code{HI} & all 32 & \code{0xFFFFFFFF} \\
\end{booktable}

Note what the bank does \emph{not} do: validate. A \code{TX_DLC} of \code{0xF} is stored and passed on
to the controller as it is \textemdash{} rejecting \code{dlc > 8} is the driver's responsibility, in
the parallel software class, and doing it twice in two languages is exactly how the two copies drift
apart. The bank masks to the bits that physically exist; the driver enforces the meaning.
\code{STATUS} and the \code{RX_*} registers ignore writes entirely, and the reserved indices 13--15
read as zero and ignore writes, exactly as the protocol specification requires.

The outputs to the controller come directly from the TX registers:
\code{tx_data <= TX_DATA_LO & TX_DATA_HI} \textemdash{} 64 bits with data byte 0 in the most
significant position, the layout \code{can_controller} expects (\chapref{c16:ch}) and the same packing
the driver performs on the other side of the map (see the data byte order diagram in
appendix~\ref{app:protocol}).

% ------------------------------------------------------------------------------------------
\section{Writing \code{register_bank.vhd}}
\label{c18:sec:write}

The file lands in \file{bridge/}, beside your copied controller modules. The shape is two processes
and a handful of concurrent assignments:
\begin{itemize}
  \item \textbf{One clocked process} owning every register and latch: the reset values (TX ready
    \code{'1'}, everything else zero), the write \code{case} over \code{reg_index}, and then the three
    updates from the controller's events, in that order.
  \item \textbf{One combinational process} for the read multiplexer: default \code{reg_rdata} to all
    zeros, and then overwrite the bits every register actually has \textemdash{} the same
    ``assign a default, override in specific branches'' style as every counter and state machine from
    \chapref{c6:ch} onwards.
  \item Concurrent assignments wiring out \code{tx_id}, \code{tx_dlc}, \code{tx_data} and
    \code{tx_req} from the TX registers and the pulse flip-flop.
\end{itemize}

Analyse it and then run the testbench handed out. \file{can_def.vhd} is a directory away, in
\file{controller/}: the bank reads \code{can_def} for its controller-facing types but belongs on the
SPI side of the protocol boundary.

\begin{shell}
cd bridge
ghdl -a --std=93 ../controller/can_def.vhd register_bank.vhd register_bank_tb.vhd
ghdl -e --std=93 register_bank_tb
ghdl -r --std=93 register_bank_tb --assert-level=error
\end{shell}

\paragraph{What the testbench checks}
\file{register_bank_tb.vhd} works through the contract in eleven cases, in order: the reset state
(bit 0 set, everything else zero); write and read back with masking on every writable register, and
the controller-facing outputs following them; \code{TX_SEND}'s one-cycle pulse and the clearing of
ready; the write discarded while a transmission is in progress; \code{tx_done} restoring ready; a
\code{TX_SEND} write of \code{0x0} doing nothing; capturing a received frame, that it holds even as
the inputs change, that it is read-only, and \code{RX_ACK}; the error latch's edge semantics,
including the write-zero rule and a clear while the level is still high; the reserved indices; and
last the two recovery paths returning TX ready after a lost arbitration, one automatic and one driven
by \code{TX_ABORT}.

Every case comment in the file says what a failed case means, so read the testbench before you write
the module \textemdash{} it is the executable form of this chapter. The last two cases are
particularly worth reading in advance: without them the node sticks after its first lost arbitration,
and since a two-node bus provokes that in normal operation it is no corner case.
